\chapter{緒論}
\label{chapter:intro}

\section{研究背景與動機}
\label{sec:background}

台灣都市環境地處副熱帶高溫夏濕氣候帶，其戶外熱舒適威脅隨全球氣候變遷加劇而日趨嚴峻。二十一世紀以來，全球氣候變遷與極端天氣事件的發生頻率與強度顯著加劇；根據聯合國跨政府氣候變化評估委員會（IPCC）的評估報告，全球地表溫度較工業革命前已大幅上升，特別是在熱帶與副熱帶氣候帶，極端高溫與持久性熱浪已逐漸轉變為常態性的氣候背景狀態。與此同時，全球都市化進程以空前速度推進，高密度的不透水鋪面、高蓄熱人工材料、密集多層建築群體，以及大量交通與空調系統釋放的人為熱，共同促成了嚴峻的都市熱島效應（Urban Heat Island, UHI），使都市邊界層之風場與輻射場發生質的改變 \citep{oke1987boundary}。台灣地理位於副熱帶與熱帶交界，屬於典型的高溫夏濕氣候，夏季高相對濕度大幅抑制人體表面的蒸散冷卻效率，使得戶外熱壓力常態性地達到不舒適乃至危險的水準 \citep{erell2011urban}。

然而，都市熱島與微氣候研究長期以來多聚焦於氣象學與環境工程視角的量化描述，對於建築師與都市設計者在實際設計現場所面對的具體困境著墨有限。當設計師嘗試以量體配置、遮蔭策略或綠化佈局回應戶外熱舒適課題時，最直接的痛點在於：任何一次幾何修改的微氣候後果，都必須等待數十分鐘至數小時的模擬運算才能得知，使得「設計—評估—修正」的迭代迴圈在方案發展的早期階段幾乎無法即時進行。這並非單純的運算效能問題，而是根本性地限制了微氣候考量能否真正融入設計思考的節奏——當回饋延遲遠超過設計師調整方案的直覺反應時間，微氣候評估便被迫退居為方案定案後的「事後驗證」角色，而非驅動幾何決策的「前置引導」角色。都市意象研究先驅 Lynch 早在其經典著作《城市意象》中即指出，市民對都市空間品質的感知，建立於可辨識、可體驗的實體環境元素之上 \citep{lynch1960image}；本研究延伸此一設計思維脈絡，主張都市戶外熱舒適同樣應被視為設計師在方案發展早期即可感知、可即時操作的空間性能維度，而非事後才能檢核的工程指標。

如何透過都市實體空間形態的量化研究與優化、開放空間的幾何優化以及永續綠化設施的空間佈局，以實質緩解夏季公共空間的熱壓力並縮短「設計—評估」迴圈的等待時間，已成為數位建築學與都市資訊學（Urban Informatics）等跨域研究社群的重要研究課題 \citep{aman2023ai, westermann2019surrogate}。

傳統的環境評估標準大多仰賴靜態、定點式的空間採樣，然而人類在都市公共空間中的空間行為在本質上是具有高度時空複雜性的移動軌跡，此一落差構成本研究第二個切入點。行人在通勤、購物、休憩漫步的過程中，持續經歷著在二維幾何空間與一維時間軸上連續交疊的動態熱暴露；學術界已指出，適當的路徑遮蔭規劃與微氣候設計能大幅改善都市開放空間的利用率與社會活動品質 \citep{coccolo2018thermal}。若戶外空間長期暴露於極端熱壓力下，將導致都市活動大幅向室內收縮，不僅引發空調能源的惡性循環，更侵奪都市居民享受健康公共生活的權利。因此，如何利用運算設計（Computational Design）工具，將行人動線上的遮蔭幾何優化與都市開放空間的形態配置進行深度聯結，是實踐氣候調適型都市設計的核心出發點。


\section{研究問題}
\label{sec:problem}

傳統高擬真度數值模擬引擎（High-fidelity Simulation Engines）難以支援即時設計迭代，這是現行數位建築設計工作流評估戶外微氣候熱舒適度時的核心運算瓶頸。計算流體力學（CFD）軟體（如 OpenFOAM、ANSYS Fluent）、全尺度微氣候熱力學模擬軟體（如 ENVI-met），以及基於 Rhino/Grasshopper 環境的 Ladybug Tools \citep{sadeghipour2013ladybug}，是學術界與實務界最常使用的評估工具，然而這些工具在應用於生成式探索或即時設計迭代時，存在著難以跨越的運算瓶頸：

\begin{enumerate}
    \item \textbf{非即時性與高時間成本：}為求解穩態納維-斯托克斯方程（Steady-state Navier-Stokes Equations）與複雜三維輻射傳熱方程，傳統 CFD 與環境模擬軟體往往需耗費數十分鐘乃至數日才能完成單一方案的模擬計算。
    \item \textbf{極度依賴高性能硬體：}網格細化之解析度直接影響模擬準確性，高精度網格計算對運算資源有極高需求。
    \item \textbf{難以整合於早期設計階段：}高延遲、難以貫通的模擬反饋機制，導致微氣候分析在實務上淪為方案完成後的事後驗證手段，難以在幾何概念設計階段提供即時的決策引導。
\end{enumerate}

上述三項限制——非即時性與高時間成本、對高性能硬體的高度依賴、以及難以整合於早期設計階段——彼此並非各自獨立的技術瑕疵，而是同一個結構性成因（模擬引擎以求解偏微分方程為核心運算邏輯，其運算複雜度與幾何精細度呈非線性增長）所導致的三種外顯症狀。這三者共同構成本研究\textbf{第一個研究缺口}：因為傳統高擬真度模擬引擎的運算延遲從根本上排除了其作為早期設計迭代工具的可能性，所以微氣候評估始終無法真正進入方案發展的前段，僅能扮演事後驗證的角色。

即使上述運算延遲問題透過某種代理模型獲得緩解，現有大量針對都市形態微氣候最佳化的學術研究，多仍聚焦於「全空間平均最佳化」（Global Spatial Average Optimization）之框架，卻缺乏結合實際行人移動軌跡的動態評估，構成本研究\textbf{第二個研究缺口}。此類方法通常將場地離散化為均質網格，並以整個研究範圍內平均 UTCI 或平均輻射溫度之極小化作為最佳化目標，存在兩項研究盲點：其一，忽視了人類步行距離之行為閾值——行人在夏季戶外極端熱壓力下的活動距離通常限制在 400 至 800 公尺（約 5 至 10 分鐘路程）內，其感受取決於路徑上各路段熱壓力之暴露歷程，而非整個場域之平均網格數值；其二，全空間平均最佳化易導致「空間品質的稀釋」——演算法為追求全局指標最小，可能將建築物與植栽稀疏佈置於整個場地，而非集中強化行人實際通行之帶狀廊道，導致優化後的方案在統計數據上表現優異，實際面對高密度人流的動線仍暴露於極端高溫下。因為現有最佳化框架未將「三維實體幾何配置」與「人行動線熱舒適品質」相互串聯比較，所以目前的研究缺乏一種能同時回應兩者的閉環運算架構。


\section{研究目的}
\label{sec:objectives}

承上節所述之兩項研究缺口，本研究提出對應之解決策略：因為傳統高擬真度模擬引擎之運算延遲難以支援即時設計迭代（第一個研究缺口），所以本研究建構低延遲之深度學習代理模型，取代傳統模擬引擎作為早期設計階段之評估工具；因為現有全空間平均最佳化方法缺乏行人動線之動態熱暴露評估（第二個研究缺口），所以本研究另行定義以人行動線為對象之熱舒適評估指標，並將兩者整合進多目標基因演算法之決策框架，使「形態幾何」與「人行動線熱舒適品質」得以在同一套運算流程中被權衡比較。具體而言，本研究之三項研究目的分別回應如下：

\noindent\textbf{研究目的一：建立低延遲全域 UTCI 預測代理模型（回應第一個研究缺口）}

因為傳統高擬真度模擬引擎之運算瓶頸源自其以求解偏微分方程為核心之運算邏輯，所以本研究不採用傳統像素型或體素型卷積神經網路（CNN）架構，改為建構異質圖神經網路（Heterogeneous Graph Neural Network）作為微氣候預測代理模型 \citep{kipf2017semi}：將都市空間中的空氣節點、建築物等實體定義為異質結構中的相異節點類型，並以邊將空間物理風場的流體聯通性與輻射場的天空視角因子嵌入其中，即時輸出網格點的 UTCI 分布。為進一步區分不同物理因果關係之訊號貢獻，本研究採用關係圖卷積網路（Relational GCN）取代同質圖卷積 \citep{schlichtkrull2018modeling}，並以真實都市建物輪廓與樓高資料取代單純程序化生成之幾何，使模型驗證得以錨定於實際都市紋理，詳見第~\ref{sec:building_extraction}節。

\noindent\textbf{研究目的二：建立熱舒適人行動線設計評估指標（回應第二個研究缺口）}

因為全空間平均最佳化無法反映行人沿路徑之動態熱暴露歷程，所以本研究以 GNN 代理模型輸出之高擬真即時空間 UTCI 數據場為基礎，建立\textbf{熱舒適人行動線設計}（Thermal Comfort-Based Walkway Design）評估指標，將行人動線網絡之熱暴露評估直接整合為多目標最佳化之適應度函數，取代僅評估單一起訖點之傳統尋路演算法（如 $A^*$）：都市開放空間設計決策之對象通常是整體人行動線系統在一整日各時段之熱舒適表現，而非單一使用者的單次行程，此一評估尺度之調整，詳見第~\ref{subsec:from_astar_to_walkway}節。

\noindent\textbf{研究目的三：整合多目標基因演算法之生成設計工作流（回應目的一、二之技術整合需求）}

因為研究目的一、二分別解決了「即時預測」與「動線評估指標」兩項獨立問題，但兩者仍須被統一至同一套設計決策流程才能實際支援方案生成，所以本研究將前兩項深度學習代理模型與熱舒適人行動線評估指標，深度整合至第二代非支配排序遺傳演算法（NSGA-II）\citep{deb2002fast} 之運算框架中，建構完整的多目標生成式決策支援系統：演算法優化多層建築物之樓層平面幾何形態組合、建築量體與高度，以及行道樹之各項參數分佈。透過多目標最佳化之迭代，系統自主找出並輸出非支配的 Pareto 最佳設計解集，供設計師在「全域熱舒適最大化」「人行動線熱舒適最大化」與「永續綠地率」之間進行理性的空間決策。

圖~\ref{fig:overview_blueprint} 呈現上述三項研究目的與後續各章節之對應關係，作為本論文之整體研究藍圖。

\begin{figure}[htbp]
    \centering
    \includegraphics[width=0.92\textwidth]{../urban-thermal-gnn/04_training/figures/fig_arch_clean.png}
    \caption{研究藍圖}
    \label{fig:overview_blueprint}
\end{figure}

上圖由基本 GIS 資訊建構（第三章）與都市形態模擬取得訓練資料，經異質圖轉譯後由 PI-ST-GNN 進行低延遲推理（第三、四章），最終驅動 NSGA-II 多目標最佳化與熱舒適人行動線設計評估，並以 Grasshopper 介面回饋予設計者（第三、五章）。


\section{研究方法概述}
\label{sec:method_overview}

本研究採用之研究方法可歸納為四個階段，完整技術細節於第三章「系統分析與設計」中依序展開：（一）以政府開放 GIS 圖資與 IoT 感測器資料建立研究場域之地理資訊基礎；（二）以 Ladybug Tools／EnergyPlus 執行都市微氣候逐時模擬，取得監督訓練所需之熱壓力標籤；（三）建構異質圖神經網路並以物理引導損失函數約束訓練，使模型輸出兼具數值精度與熱力學一致性；（四）將訓練完成之代理模型嵌入 NSGA-II 多目標最佳化與熱舒適人行動線設計評估，並透過 FastAPI／WebSocket \citep{rfc6455} 與 Rhino Grasshopper 介面整合，形成可即時互動之閉環設計決策系統。

上述四個階段與前節三項研究目的之對應關係如下：階段（一）至（三）——地理資訊建構、微氣候模擬取樣、異質圖神經網路訓練——共同支撐\textbf{研究目的一}之低延遲全域 UTCI 預測代理模型；階段（四）中之熱舒適人行動線評估機制對應\textbf{研究目的二}，NSGA-II 多目標最佳化與 Grasshopper 介面整合則對應\textbf{研究目的三}。此一方法論安排之邏輯，在於使每一階段之技術選擇均直接回應第~\ref{sec:problem}節所述之研究缺口，而非孤立的技術堆疊。


\section{研究範圍與限制}
\label{sec:scope}

本研究在空間尺度的假定上，嚴格聚焦於都市「行區尺度」（Neighborhood Scale / District Scale），具體運算幾何規模範圍設定在邊長約 \textbf{500 公尺至 800 公尺}的正方形都市行區範疇。此一尺度的選定具有充足的科學與實務依據：

\begin{enumerate}
    \item \textbf{環境學依據：}根據都市微氣候學奠基者 Oke 的理論，行區尺度是都市微氣候中「微氣候混合層」與「街道峽谷（Urban Canyons）」次系統的核心交點。在此尺度上，都市表面幾何形態（如建築高度豐富質性、天空視角因子 SVF）與局部風場的空氣動力學粗糙長度（Roughness Length）產生最強的交互作用，是直接決定 pedestrian-scale（人體尺度）微氣候特徵的關鍵尺度 \citep{oke1987boundary}。
    \item \textbf{運算限制：}儘管本研究採用了高效的 GNN 代理模型，但下游 NSGA-II 多目標最佳化之族群個體迭代，仍受到節點數量之指數增長運算負載限制。若下至單棟建築實體尺度，則無法覆蓋都市行區風廓與群體遮蔭的巨觀行為；若上擴至都市總體規劃尺度，則將超出當前計算幾何在 Rhino/Grasshopper 前端資料傳輸 Payload 的可能極限。因此，500m--800m 的行區範疇是兼顧環境行為學實用性與計算幾何可行性的最佳尺度。
\end{enumerate}

在時間軸與邊界條件的設定條件上，本研究聚焦以夏季夏至熱浪日（6 月 21 日前後）之日間時段（08:00 至 18:00）。此一時間限制主要依據台灣微氣候學研究針對熱帶與副熱帶高密度都市空間之熱舒適研究結論：台灣高溫夏濕環境中的最高熱壓力與都市熱島強度峰值，高度集中於夏至日前後正午至下午之間，此時段的太陽高度角最大，直接輻射與地表二次長波輻射釋放已達飽和，戶外 UTCI 常態性突破 $38\,^\circ\text{C}$，達到強烈乃至極端的熱壓力風險，且為都市行人暴露於紫外線與高溫危害、引發急性公共健康風險的關鍵時段 \citep{su2023outdoor}。以該時段的微氣候邊界條件作為 GNN 代理模型的靜態氣象特徵，對於評估極端環境適應性、建立高防禦性的生成式設計規範，具有統計代表性與極端工程設計（Worst-case Scenario Engineering）的科學效度。
